% Exported from Math++. Compile with pdflatex.
\documentclass[11pt,a4paper]{article}
\usepackage[T1]{fontenc}
\usepackage[utf8]{inputenc}
\usepackage{lmodern}
\usepackage{amsmath,amssymb,amsthm,mathtools}
\usepackage{graphicx,booktabs,xcolor}
\usepackage[hidelinks]{hyperref}
\usepackage[margin=25mm]{geometry}
\providecommand{\degree}{\ensuremath{^\circ}}
\providecommand{\xmark}{\ensuremath{\times}}
\providecommand{\omicron}{o}
\providecommand{\oiint}{\mathop{\oint\mkern-12mu\oint}}
\theoremstyle{plain}
\newtheorem{theorem}{Theorem}
\newtheorem{lemma}[theorem]{Lemma}
\newtheorem{corollary}[theorem]{Corollary}
\newtheorem{proposition}[theorem]{Proposition}
\newtheorem{conjecture}[theorem]{Conjecture}
\theoremstyle{definition}
\newtheorem{definition}[theorem]{Definition}
\newtheorem{example}[theorem]{Example}
\theoremstyle{remark}
\newtheorem*{remark}{Remark}
\setlength{\parindent}{0pt}\setlength{\parskip}{0.6em}
\title{Show the working}
\date{}
\begin{document}
\maketitle
Write \texttt{= ?!} instead of \texttt{= ?} and every answer comes with its working, one rule per step.
Turn on "Show steps" to get the working for every \texttt{= ?} query.

\section*{Derivatives}

The product, chain and power rules, each named as it is used.

\begin{equation*}
  \frac{d}{dx} x^2 \sin(3x) = 3 x^{2} \cos\left(3 x\right) + 2 x \sin\left(3 x\right)
\end{equation*}
\begin{equation*}
  \frac{d}{dx} \frac{\sin x}{x^2} = \frac{x^{2} \cos\left(x\right) - 2 x \sin\left(x\right)}{x^{4}}
\end{equation*}
\section*{Integrals}

A substitution writes out $u$ and $du$:

\begin{equation*}
  \int 2x \cos(x^2) \, dx = \sin\left(x^{2}\right) + C
\end{equation*}
Integration by parts names $u$, $dv$, $du$ and $v$:

\begin{equation*}
  \int x e^x \, dx = e^{x} x - e^{x} + C
\end{equation*}
Partial fractions show the decomposition:

\begin{equation*}
  \int \frac{1}{x^2 - 1} \, dx = \frac{\ln\left(\left|x - 1\right|\right)}{2} - \frac{\ln\left(\left|x + 1\right|\right)}{2} + C
\end{equation*}
A definite integral ends with the fundamental theorem of calculus:

\begin{equation*}
  \int_0^2 x^2 \, dx = \frac{8}{3}
\end{equation*}
\section*{Limits}

L'H\^{o}pital's rule says which indeterminate form it removes:

\begin{equation*}
  \lim_{x \to 0} \frac{e^x - 1 - x}{x^2} = \frac{1}{2}
\end{equation*}
\section*{Series}

The Taylor coefficients are listed:

\begin{equation*}
  \operatorname{Taylor\ series}_{x=0}^{4} e^x \sin x = x + x^{2} + \frac{x^{3}}{3} + O\left(x^{4}\right)
\end{equation*}
\end{document}
